\documentclass[twocolumn]{aastex631}
\begin{document}
\title{Mass-metallicity relation}
\author{NebulaMind Lab (autonomous pipeline)}
\affiliation{NebulaMind Lab, \url{https://lab.nebulamind.net}}
\begin{abstract}
Mass-metallicity relation — median relations for SDSS (120,000 gals). Generated autonomously from public data (sdss) via the NebulaMind Lab runner: a bounded, reproducible, descriptive study.
\end{abstract}
\section{Introduction}
The study of cosmic-chemical-evolution is crucial for understanding how galaxies form and evolve over time. One key aspect of this evolution is the relationship between a galaxy's mass and its metallicity, which provides insights into the processes that govern star formation and gas exchange within galaxies. By examining this mass-metallicity relation, we can gain a better understanding of the underlying mechanisms driving cosmic-chemical-evolution.
\section{Data and method}
To investigate the mass-metallicity relation, we utilized data from the Sloan Digital Sky Survey (SDSS), which provided a vast dataset containing information on approximately 120,000 galaxies. We employed a method that focused on analyzing the median relations between galaxy mass and metallicity within this extensive sample. This approach allowed us to identify trends and patterns in the data without relying on complex models or assumptions.
\section{Result}
Our analysis revealed a clear mass-metallicity relation within the SDSS dataset. Specifically, we found that there are median relations for the 120,000 galaxies studied, indicating a systematic relationship between galaxy mass and metallicity. This result supports the notion that more massive galaxies tend to have higher metallicities, providing valuable insights into the processes governing cosmic-chemical-evolution.
\begin{figure}\centering\includegraphics[width=\columnwidth]{result.png}\caption{Mass-metallicity relation}\end{figure}
\section{Caveats}
However, it is essential to acknowledge the limitations of our study. The use of an automated, single-selection method may introduce biases or oversights in our analysis, as it relies on a specific set of criteria that might not capture the full complexity of galaxy properties. Additionally, our measurements were uncalibrated, which means they have not been adjusted for potential systematic errors or compared to other independent datasets. These factors could impact the accuracy and reliability of our findings, highlighting the need for further research and validation using more comprehensive and refined methods.
\end{document}
